\documentclass[11pt,a4paper]{article}
\usepackage[T1]{fontenc}
\usepackage[utf8]{inputenc}
\usepackage{amsmath,amssymb,xcolor}
\usepackage[margin=25mm]{geometry}
\usepackage[colorlinks=true,urlcolor=blue]{hyperref}
\definecolor{nexorablue}{HTML}{2459B5}
\setlength{\parindent}{0pt}
\setlength{\parskip}{7pt}
\setlength{\emergencystretch}{2em}
\renewcommand{\familydefault}{\sfdefault}
\begin{document}

{\small\bfseries\color{nexorablue} NEXORA CPGE}\par
{\small MP / NEX-IND-04 / Corrigé}\par
{\LARGE\bfseries Puissance nécessaire pour tirer une barre sur des rails}\par
\vspace{0.25em}\hrule\vspace{0.6em}

\section*{Électricité}
Les valeurs absolues sont $|e|=B\ell v=0{,}28\,\mathrm V$ et $|i|=0{,}35\,\mathrm A$.

\section*{Mécanique et bilan}
La force magnétique vaut $B\ell|i|\simeq0{,}039\,\mathrm N$ en sens opposé au déplacement. Il faut exercer la force opposée. Sa puissance vaut $0{,}098\,\mathrm W$, comme $Ri^2$.


\end{document}
